\chapter{L'organisation de l'équipe}
\label{ch:organisation}

\section{Les principes}

\begin{itemize}
  \item \textbf{Un module, un responsable, un suppléant.} Le responsable décide dans son module ;
        le suppléant connaît le code assez pour le reprendre si besoin, et relit chaque fusion.
  \item \textbf{Des contrats figés lundi soir.} Les modules communiquent par des structures de
        données définies ensemble, une fois. Chacun développe contre des \emph{bouchons} qui
        respectent le contrat : personne n'attend personne.
  \item \textbf{\f{main} tourne toujours.} À tout moment, \f{streamlit run app.py} doit lancer la
        démo depuis \f{main}.
  \item \textbf{Deux intégrations mardi}, à 12h30 et à 18h, avec le jeu de scénarios comme juge.
  \item \textbf{D'abord une solution sûre, ensuite les idées.} Jusqu'à J2, on construit ce qui est
        prévu. Après, on essaie les idées nouvelles, même exotiques : c'est ce que l'organisateur
        nous a encouragés à faire, une fois la base solide. Et on se fait plaisir.
\end{itemize}

\section{Les trois modules}

Le découpage suit la cascade, pas les filières. Chaque module contient une part de données, une
part de modèle et une part de logiciel : chacun touche à l'IA, au code et à la mesure, et chacun
peut présenter son module au jury avec des chiffres.

\begin{center}\small
\begin{tabularx}{\textwidth}{@{}L{3.3cm} Y Y Y@{}}
\toprule
& \textbf{Module 1 --- Comprendre et mesurer} & \textbf{Module 2 --- Converser et formuler} &
  \textbf{Module 3 --- Conseiller et agir} \\
\midrule
\textbf{Responsable} & Younes & Yahya & Sohayb \\
\textbf{Suppléant} & Yahya & Sohayb & Younes \\
\textbf{Étages} & 1, et la mesure transversale & 2, 5, l'escalade, l'interface & 3 et 4 \\
\textbf{Part données} & Banking77 en français, exemples des 8 intentions & Glossaire, gabarits &
  Clients simulés, catalogue \\
\textbf{Part modèle} & Mots-clés, encodeur, têtes, seuil $\tau$ & Choix du petit LLM, garde-fou & Modèle
  d'appétence \\
\textbf{Part logiciel} & \f{comprendre()}, instrumentation, bancs de mesure & \f{reply()},
  automate de dialogue, Streamlit & Moteur de règles, outils, rejeu des scénarios \\
\textbf{Dans le pitch} & Les résultats et la frugalité & Le problème, l'architecture, la démo &
  L'explicabilité et la conformité \\
\bottomrule
\end{tabularx}
\end{center}

\subsection{Module 1 --- Comprendre et mesurer (Younes)}

Le module où se trouve le calcul lourd --- inférence par lots sur DeepEirb, bancs de mesure
d'énergie --- et aussi un vrai travail de modélisation.

\begin{itemize}
  \item Obtenir l'accès à DeepEirb ; traduire Banking77 (vLLM ou, à défaut, Marian sur la 3060) ;
        générer et faire relire les exemples des 8 nouvelles intentions ; écrire la fiche de
        données.
  \item Construire le filtre par mots-clés et sa règle de confiance ; entraîner l'encodeur et les
        deux têtes (famille, intention) ; choisir $\tau$ par balayage ;
        écrire les règles d'extraction des entités ; exposer \f{comprendre(message, etat)}.
  \item Remplir le tableau de compréhension du chapitre~\ref{ch:evaluation}.
  \item Écrire l'instrumentation (\f{src/mesure/frugalite.py}) et le banc A/B/C ; produire le
        graphique du pitch.
  \item En bonus : faire tourner l'étage 1 et les règles sur la Jetson.
\end{itemize}
\textbf{Fichiers} : \f{data/processed/banking77\_fr/}, \f{data/processed/intentions\_fr.csv},
\f{src/comprendre/}, \f{src/mesure/}, \f{eval/bench\_nlu.py}, \f{eval/bench\_frugalite.py}.

\subsection{Module 2 --- Converser et formuler (Yahya)}

Le module de l'orchestration : il appelle les deux autres et porte l'interface.

\begin{itemize}
  \item Écrire l'automate de dialogue et les neuf parcours YAML ; écrire \f{reply()}, le point
        d'entrée unique de la cascade.
  \item Écrire les gabarits ; brancher le petit LLM via Ollama ; écrire le garde-fou ; dérouler le
        protocole de choix du modèle de formulation.
  \item Brancher l'escalade (DeepEirb ou Ollama) ; écrire le glossaire et sa recherche en mémoire.
  \item Construire l'interface : conversation, profil, panneaux \og pourquoi \fg{} et \og
        frugalité \fg, mode mots simples, messages de démo ; enregistrer la vidéo de secours.
\end{itemize}
\textbf{Fichiers} : \f{src/assistant.py}, \f{src/converser/}, \f{app.py},
\f{data/processed/glossaire.json}.

\subsection{Module 3 --- Conseiller et agir (Sohayb)}

Le module des décisions : il doit être juste, testé et lisible, ce qui en fait un module de génie
logiciel autant que d'IA.

\begin{itemize}
  \item Écrire l'adaptateur de la base clients et l'enrichissement simulé ; créer les huit clients
        de démonstration.
  \item Écrire le catalogue JSON avec sources et dates ; écrire le moteur de règles (éligibilité,
        pertinence, classement, raisons, produits écartés) et les calculs de crédit.
  \item Entraîner et évaluer le modèle d'appétence ; exposer ses trois facteurs principaux.
  \item Écrire les cinq outils, la confirmation et le journal.
  \item Écrire le programme de rejeu des scénarios ; le lancer à chaque intégration ; tenir le dépôt
        Git propre jusqu'au dépôt final.
\end{itemize}
\textbf{Fichiers} : \f{src/conseiller/}, \f{data/processed/clients.csv},
\f{data/processed/catalogue.json}, \f{eval/run\_scenarios.py}, \f{tests/}.

\subsection{Ce qui est commun}

Chacun écrit les scénarios de son module (au moins dix). Chacun écrit les diapositives de son
module. Les contrats sont écrits ensemble, lundi soir, et ne changent plus qu'à l'unanimité.

\section{Les contrats entre modules}

Un seul fichier, \f{src/contrats.py}, importé par tous. Les noms de champs sont en français et sans
accents.

\begin{code}
\footnotesize
\begin{verbatim}
from dataclasses import dataclass, field

@dataclass
class Comprehension:                  # rendu par le module 1
    famille: str                      # "credit", "carte_perdue", ...
    intention: str                    # "credit_immobilier", "lost_or_stolen_card", ...
    confiance: float                  # probabilite de la famille, entre 0 et 1
    alternatives: list                # [("virement", 0.21), ...] pour la clarification
    entites: dict                     # {"montant": 250000, "lien": "fils", "objet": "voiture"}
    insatisfaction: bool = False
    modele: str = "e5-small+lr"       # ou "mots-cles", ou "ministral-3:14b" (escalade)
    exemple_proche: str = ""          # phrase d'entrainement la plus proche (explication)

@dataclass
class Proposition:                    # rendu par le module 3
    produit_id: str                   # "livret_a"
    raisons: list                     # ["Lucas a 20 ans : Livret Jeune possible (12-25 ans)"]
    sources: list                     # URL publiques du catalogue
    score: float = 0.0

@dataclass
class Conseil:                        # rendu par le module 3
    propositions: list                # [Proposition], triees
    ecartes: list                     # [("lep", "revenu au-dessus du plafond")]
    manque: list                      # ["prix_bien"] : informations a demander
    alertes: list                     # ["taux d'effort 42,1 % > 35 %"]
    faits: dict                       # tous les nombres que le texte a le droit de citer

@dataclass
class Action:                         # rendu par le module 3
    outil: str                        # "bloquer_carte"
    parametres: dict
    confirmee: bool
    resultat: dict = field(default_factory=dict)   # {"reference": "OPP-..."}

@dataclass
class Etape:                          # une ligne du panneau frugalite (module 1)
    nom: str; modele: str; latence_ms: float; tokens: int = 0

@dataclass
class Reponse:                        # rendu par le module 2, affiche par l'interface
    texte: str
    comprehension: Comprehension
    conseil: Conseil | None
    action: Action | None
    escalade: bool
    trace: list                       # [Etape]

def reply(message: str, client_id: str, etat: dict) -> tuple[Reponse, dict]: ...
\end{verbatim}
\end{code}

\section{L'arborescence du dépôt}

\begin{code}
\footnotesize
\begin{verbatim}
hackathon-IA/
|-- app.py                       interface Streamlit                      (M2)
|-- src/
|   |-- contrats.py              structures partagees                     (tous)
|   |-- assistant.py             reply() : la cascade                     (M2)
|   |-- comprendre/              nlu.py, entites.py, tetes/               (M1)
|   |-- converser/               dialogue.py, parcours/*.yaml,
|   |                            formuler.py, garde_fou.py, glossaire.py  (M2)
|   |-- conseiller/              clients.py, catalogue.py, regles.py,
|   |                            appetence.py, outils.py                  (M3)
|   `-- mesure/                  frugalite.py                             (M1)
|-- data/
|   |-- raw/                     Banking77, Bank Marketing (ignore par git)
|   `-- processed/               jeux francais, clients, catalogue, glossaire
|-- eval/
|   |-- scenarios/*.yaml         30 conversations                         (tous)
|   |-- run_scenarios.py                                                  (M3)
|   |-- bench_nlu.py, bench_frugalite.py                                  (M1)
|-- tests/                       tests unitaires des regles               (M3)
`-- rapport/                     ce document
\end{verbatim}
\end{code}

Les données téléchargées sont aujourd'hui dans \f{data/banking77/} et \f{data/bank-marketing/} :
elles passent dans \f{data/raw/}, ajouté au \f{.gitignore}. Les jeux préparés, petits et produits
par nous, sont versionnés : ils font partie du livrable.

\section{Le planning}

\begin{table}[H]
\centering\scriptsize
\renewcommand{\arraystretch}{1.25}
\begin{tabularx}{\textwidth}{@{}L{1.45cm} Y Y Y Y Y@{}}
\toprule
& \textbf{Lundi\newline 18h--20h} & \textbf{Mardi\newline 8h30--12h30} & \textbf{Mardi\newline 13h30--18h}
  & \textbf{Mardi\newline 18h--20h (option)} & \textbf{Mercredi\newline 8h30--12h} \\
\midrule
\textbf{Younes}\newline (M1) & \cellcolor{ent}Accès DeepEirb ; lancer la traduction
  & \cellcolor{ent}Mots-clés ; têtes v1 (anglais et français partiel) ; entités ; \texttt{comprendre()}
  & \cellcolor{ent}Français complet ; 8 intentions ; seuil $\tau$ ; escalade ; instrumentation
  & \cellcolor{ent}Banc A/B/C sans surveillance
  & \cellcolor{ent}Chiffres finaux ; graphique ; Jetson (bonus) \\
\textbf{Yahya}\newline (M2) & \cellcolor{tab}Parcours carte ; squelette de \texttt{app.py}
  & \cellcolor{tab}Automate ; parcours carte et frais ; gabarits ; petit LLM branché
  & \cellcolor{tab}Parcours épargne, crédits, virement, insatisfaction ; garde-fou ; panneaux
  & \cellcolor{tab}Scénarios en échec
  & \cellcolor{tab}Vidéo de secours ; finitions de l'interface \\
\textbf{Sohayb}\newline (M3) & \cellcolor{ext}8 clients ; catalogue v0
  & \cellcolor{ext}\texttt{clients.py} ; règles épargne et crédit ; opposition ; frais
  & \cellcolor{ext}Appétence ; produits écartés ; virement ; conseiller ; rejeu et scénarios
  & \cellcolor{ext}Scénarios en échec
  & \cellcolor{ext}Dépôt Git final ; README \\
\textbf{Tous} & \cellcolor{mod}Lire ce document ; contrats figés ; bouchons
  & & & & \cellcolor{mod}Diapositives ; deux répétitions ; dépôt à 11h30 \\
\midrule
\textbf{Jalon} & \textcolor{kored}{\textbf{J0}} à 20h & \textcolor{kored}{\textbf{J1}} à 12h30 :
  intégration 1 & \textcolor{kored}{\textbf{J2}} à 18h : intégration 2, gel & ---
  & \textcolor{kored}{\textbf{J3}} à 11h30 : tout est remis \\
\bottomrule
\end{tabularx}
\caption{Le planning, par demi-journée.}
\label{tab:planning}
\end{table}

\begin{center}\small
\begin{tabularx}{\textwidth}{@{}l l Y@{}}
\toprule
\textbf{Jalon} & \textbf{Quand} & \textbf{Critère de passage} \\
\midrule
J0 & Lundi, 20h & \f{src/contrats.py} écrit et fusionné ; \f{streamlit run app.py} répond avec
  les bouchons ; la traduction est lancée (ou son repli décidé). \\
J1 & Mardi, 12h30 & Intégration 1 : Inès (carte) et Marc (frais) passent de bout en bout, avec
  les vrais modules. \\
J2 & Mardi, 18h & Intégration 2 : les neuf parcours passent leurs scénarios principaux ; la trace
  est complète. \textbf{Gel fonctionnel} : après J2, on corrige et on mesure, on n'ajoute plus. \\
J3 & Mercredi, 11h30 & Tableaux de résultats remplis ; vidéo enregistrée ; deux répétitions
  chronométrées faites ; code déposé (étiquette \f{v1.0}) ; support remis. \\
\bottomrule
\end{tabularx}
\end{center}

\section{Les priorités}

\begin{center}\small
\begin{tabularx}{\textwidth}{@{}L{2.6cm} Y@{}}
\toprule
\textbf{Priorité} & \textbf{Contenu} \\
\midrule
\textbf{Indispensable} (J2) & Étage 1 mesuré en anglais et en français ; parcours carte, frais,
  épargne, appartement, crédit auto ; règles avec raisons et produits écartés ; les 8 clients ;
  petit LLM et garde-fou ; panneaux pourquoi et frugalité ; banc A/B/C. \\
\textbf{Important} & Virement en mots simples ; insatisfaction et passage de main ; jargon ; prêt
  étudiant ; appétence ; escalade. \\
\textbf{Si le temps le permet} & Idées nouvelles, même exotiques, encouragées par l'organisateur une
  fois la base sûre ; déploiement sur Jetson ; réponse dans la langue du client ;
  outils exposés en MCP ; saisie vocale locale (Whisper), utile à l'accessibilité. \\
\textbf{Hors périmètre} & Affinage d'un LLM ; connexion à un vrai système bancaire ;
  authentification ; base de données serveur. \\
\bottomrule
\end{tabularx}
\end{center}

\section{Les règles de travail}

\begin{itemize}
  \item Branches courtes \f{m1/<sujet>}, \f{m2/<sujet>}, \f{m3/<sujet>}, fusionnées dans la
        demi-journée, après relecture du suppléant et passage des scénarios.
  \item Aucune donnée brute ni aucun secret dans Git ; les modèles ne sont jamais versionnés.
  \item Un point de cinq minutes toutes les trois heures : fait, bloqué, suivant.
  \item Tout résultat chiffré est noté au moment où il est obtenu, avec la commande qui l'a
        produit : le support se construit au fil de l'eau.
  \item \textbf{Définition de \og fini \fg{} pour un module} : le code est dans \f{main}, ses
        scénarios passent, ses mesures sont notées, et sa diapositive est écrite.
\end{itemize}
